\begin{exercise}{112}
Let $P$ be a poset and $\mathcal C$ an $\infty$-category. Suppose that you are given a function
$f\colon P\to\mathrm{ob}(\mathcal C)$ with the following property: Whenever $x\le y$ are elements
of $P$, then the spaces of maps $\mathrm{map}_{\mathcal C}(fx,fy)$ are contractible.
\begin{enumerate}[label=(\arabic*)]
 \item Show that there exists a functor $F\colon P\to\mathcal C$ extending the given function $f$
       on objects.
 \item Show that any two such extensions are equivalent.
\end{enumerate}
Of course, the ``correct'' version of (2) is the following: Show that there is a contractible space
parametrizing all possible choices of such extensions, i.e., that in the pullback diagram
\[
\begin{tikzcd}[column sep=large,row sep=large]
 \mathrm{Ext}(f)\arrow[r]\arrow[d] & \mathrm{Fun}(P,\mathcal C)\arrow[d] \\
 \Delta^0\arrow[r] & \mathrm{Fun}(P_0,\mathcal C)
\end{tikzcd}
\]
the simplicial set $\mathrm{Ext}(f)$ is a contractible Kan complex. Can one replace $P$ by an arbitrary
$1$-category and obtain the same results?

\begin{proof}
We work with an arbitrary ordinary category $D$.
Assume that $\mathrm{map}_{\mathcal C}(fa,fb)$ is contractible whenever
there is a morphism $a\to b$ in $D$.
Write $K=N(D)$ and let $V\subseteq K$ be its discrete set of vertices.
Let $q\colon\mathcal C\to N(h\mathcal C)$ be the canonical map.
Whenever $D(a,b)$ is nonempty, we have
$\mathrm{Hom}_{h\mathcal C}(fa,fb)=\{*\}$.
Indeed, every homotopy-category morphism has an edge representative
(Corollary~1.2.11), and a path in the contractible mapping space gives
squares with identity vertical edges, identifying these representatives
in $h\mathcal C$.
Thus $f$ determines a functor $\bar f\colon D\to h\mathcal C$:
each arrow goes to the unique possible morphism.
Identities and compositions are preserved by this uniqueness.

Form the pullback
\[
\begin{tikzcd}[column sep=large,row sep=large]
 E=K\times_{N(h\mathcal C)}\mathcal C\arrow[r]\arrow[d,"p"']
   & \mathcal C\arrow[d,"q"] \\
 K\arrow[r,"N\bar f"'] & N(h\mathcal C).
\end{tikzcd}
\]
The map $q$ is an inner fibration: fill an inner horn in $\mathcal C$;
its image is the prescribed simplex because an ordinary nerve has
unique inner-horn fillers (Theorem~1.1.52).
Hence $p$ is an inner fibration and $E$ is an $\infty$-category.
Its vertices are exactly $\widetilde a=(a,fa)$, so $p$ is bijective on vertices.
Exponentials preserve pullbacks, so Definition~1.3.47 gives
\[
 \mathrm{map}_E(\widetilde a,\widetilde b)
 \cong D(a,b)\times_{\mathrm{Hom}_{h\mathcal C}(fa,fb)}
                         \mathrm{map}_{\mathcal C}(fa,fb),
\]
where ordinary hom-sets are discrete simplicial sets.
If $D(a,b)$ is empty this is empty; otherwise it is a coproduct,
indexed by $D(a,b)$, of copies of the contractible mapping space.
The projection to $D(a,b)$ is therefore a homotopy equivalence.
Thus $p$ is fully faithful, and Theorem~2.3.20 makes it a Joyal equivalence.
By Exercise~\ref{ex:104}, $p$ is a trivial fibration, hence lifts against
every monomorphism by Exercise~\ref{ex:59}.

We now check the actual strict extension space.
Let $T\hookrightarrow U$ be any monomorphism and suppose
$T\to\mathrm{Ext}_D(f)$ is given.
Its adjoint $H\colon K\times T\to\mathcal C$ is constant at $f(a)$ on
every $\{a\}\times T$.
Moreover, $qH=N\bar f\operatorname{pr}_K$:
on an edge projecting to $a\to b$, both maps take the unique value
in $\mathrm{Hom}_{h\mathcal C}(fa,fb)$, and maps into an ordinary nerve
are determined by their edges and vertices.
Thus $H$ lifts to $E$ over $K$.
Together with the prescribed constant maps on $V\times U$, it defines
a map over $K$ on
\[
 (K\times T)\cup(V\times U)\hookrightarrow K\times U.
\]
The maps agree on $V\times T$.
This inclusion is a monomorphism, so $p$ supplies an extension.
Composing with $E\to\mathcal C$ gives a map
$U\to\mathrm{Ext}_D(f)$, with every vertex family still strictly constant.

Consequently $\mathrm{Ext}_D(f)\to\Delta^0$ lifts against all monomorphisms.
It is nonempty and Kan. After choosing a vertex $e$, extend the map
$\mathrm{Ext}_D(f)\times\partial\Delta^1\to\mathrm{Ext}_D(f)$
given by the identity and the constant map at $e$ across
$\mathrm{Ext}_D(f)\times\Delta^1$.
This is a contraction.
Nonemptiness gives an extension, and any two extensions are joined
by an edge of this Kan complex. Such an edge remains an equivalence
in $\mathrm{Fun}(K,\mathcal C)$.
This proves (1), (2), and the stronger assertion, and answers the final
question affirmatively for every ordinary category.
\end{proof}
\end{exercise}
